\begin{afterword}
\summaryhead

The post starts from allegations in The Register that Wikipedia's senior administrators used a secret mailing list to ban critics, once banning a productive editor as a supposed spy. It rejects the idea that something in the cause of organizing the world's knowledge produced this. Ordinary human nature, meaning in-group bias and the spirals of praise and hate described in earlier posts, is enough.

As heat evens out and programs gather patches, the author argues, every cause ``wants'' to be a cult. Any group with an unusual goal drifts toward cultishness unless it works constantly against the drift, and calling the cause ``rationality'' does not exempt it, as the Objectivists showed. Cultishness is a matter of degree, even in science, where one can ask whether journals favor famous authors and whether blind review helps. Admitting such weak points does not betray a cause. Believing that only foolish other causes turn into cults is what stops a group from making the effort.

\responsehead

The post applies its warning to its own side. It names a group whose cry was ``Rationality! Reason! Objective reality!'' and says the label protects nothing. That is to its credit, and it leads to ``Two Cult Koans,'' nine days later.

The central claim covers every cause, and its support is an analogy. ``Every Cause,'' ``Every group of people with an unusual goal'': the evidence is one set of allegations about Wikipedia and a footnoted reference to the Objectivists. The rest is comparison with heat and with software, which can illustrate a drift but cannot show that human groups have one. Calling cultishness a ``high-entropy state'' and an ``attractor'' makes the claim sound like a law of physics without adding any evidence for it.

The claim also has no stated test. Cultishness is said to be ``quantitative,'' but no quantity is offered. Resisting it takes ``constant effort,'' but the post does not say what counts as the effort, so it names no case that could count against the claim.

The one measurable example, bias in journal review, is posed as a list of questions. Evidence on the first question had been published a quarter-century earlier: resubmitted under unknown names and institutions, eight of nine undetected published articles were rejected. Citing it would have given the essay the kind of measurement its own view of cultishness as ``quantitative'' calls for.

In short: a warning the author applies to the author's own side, supported by analogy and one set of allegations, and stated without a measure or a test.
\end{afterword}
